\chapter{Generators}
\label{ch:generators}

\section{Water Wheel}

Treated wood. Mounted on a \textbf{Dynamo}, turned by flowing water down one side. Output scales
with how much of the wheel is being driven, and three wheels can share one dynamo.

Cheap, silent, needs no fuel, produces very little. It is the right first generator and never the
last.

\section{Windmill}

Treated wood on a dynamo again, mounted high and in the open. Output rises with altitude and falls
if obstructed. The \textbf{Improved Windmill} adds iron sheet metal for a substantial increase.

\begin{ietip}
Windmills need clear air around them, not merely above them. A windmill in a valley is a decoration.
\end{ietip}

\section{Thermoelectric Generator}

Two blocks of differing temperature on opposite sides --- classically lava and packed ice, or a
blaze-adjacent block against a cold one --- and it produces a steady trickle with no moving parts
and no fuel logistics. The default output multiplier is~1.

Its virtue is that it never stops. Its vice is that it is slow.

\section{Diesel Generator}

The first serious generator: a multiblock burning biodiesel for \textbf{4\,096 IF/t} by default.
Needs a fluid supply, and rewards a real one.

\begin{iefork}
The petroleum chain supplies it far better than the biodiesel loop does, and adds several larger
plants above it, ending at the Gas Turbine. \autoref{ch:petroleum}.
\end{iefork}

\section{Lightning Rod}

A tall multiblock that captures strikes for \textbf{16\,000\,000 IF} apiece by default. Entirely
weather-dependent, spectacular, and useless as a base load --- pair it with storage large enough to
hold a strike or most of it is wasted.

\section{Keeping a Distant Plant Running}
\label{sec:generation-meter}

A generator only works while its chunk is loaded. Wire carries power through unloaded chunks, but a
plant nobody is standing near produces nothing, so the town at the far end of the line goes dark the
moment you leave the plant behind.

\begin{iefork}
The \textbf{Generation Meter} solves this without a chunk loader. Place it against the generator's
output face and bolt the wire connector onto its front. Power passes straight through; the meter
stores none of its own and measures what the generator produces. Right-click it with an empty hand for a readout.

While the plant's chunks are unloaded, the wire network keeps receiving that measured output at every
consumer that is still loaded. Nothing is kept loaded to do it.
\end{iefork}

\begin{description}
\item[Let it measure] Leave a new meter beside the running plant for a minute. It reports the best
      output it saw over that minute, not the last reading.
\item[Fuel] Outside city mode only fuel-free generators keep supplying while unloaded --- a dynamo on
      a water wheel or windmill, or a thermoelectric generator. In city mode every generator does, and
      burns no fuel while unloaded.
\item[Consumers] Still have to be loaded. The meter keeps the plant running, not the town.
\item[Grid Feed Units] A meter in front of a Feed Unit instead of a connector keeps that unit supplying
      its segment while unloaded; the console lists it as \emph{virtual}. \autoref{ch:grid}.
\end{description}

\begin{ietip}
Operators can list every metered plant, and whether it is running for real or virtually, with
\texttt{/ie virtualgen}.
\end{ietip}

\section{Storage}

\begin{description}
\item[LV / MV / HV Capacitor] Buffers of increasing size. Each face can be configured to input,
      output or nothing with a hammer. As placed, only the \emph{top} face accepts input and every
      other face is set to nothing --- a fresh capacitor fed from the side takes no charge at all
      until you configure that face.
\item[Creative Capacitor] A bottomless supply for testing, outputting on all six faces.
\item[Capacitor Backup] Keeps a capacitor's contents through a chunk unload.
\end{description}

\begin{iewarning}
A capacitor with every face set to output will happily drain itself into whatever is adjacent.
Configure the faces deliberately.
\end{iewarning}

\begin{iefork}
City mode makes fuel cosmetic for generators: a presence check and a token sip instead of a per-tick
burn rate and tank drain. A generator with fuel in it runs; how much is left stops being tracked.
\end{iefork}
